\begin{textsample}{POS Dim 4 – ro – Score 9.00 – ro/ro\_ef\_30.txt}  \label{ex:f4_pos_218}
7.3 EDUCAÇÃO ESCOLAR \textbf{INDÍGENA} Os movimentos \textbf{indígenas} e as autoridades competentes em todas as esferas de atuação política possibilitaram a formalização de leis federais e estaduais, além da criação de pareceres e resoluções que viabilizaram a legalização da Educação Escolar \textbf{Indígena} específica, diferenciada, \textbf{bilíngue} e intercultural.
Dentre a legislação que norteia a educação escolar \textbf{indígena} no Brasil destacam -se : Constituição Federal de 1988 ; Decreto Presidencial, No 26/91 ; Lei de Diretrizes e Bases da Educação Nacional ; Plano Nacional de Educação ; Parecer 14/99 ; Resolução 03/99.
A Educação Escolar \textbf{Indígena} no Brasil vem obtendo, desde a década de 70, avanços significativos no que diz respeito à legislação que a regula.
Existem hoje leis bastante favoráveis quanto ao reconhecimento da necessidade de uma educação específica, diferenciada e de qualidade para as \textbf{populações} \textbf{indígenas}.
Hoje, segue -se a orientação dos pressupostos sugeridos nos manuais disponibilizados pelo Governo Federal, como parâmetro para um desenvolvimento escolar alinhado aos saberes e processos próprios de aprendizagem específicos de cada povo, levando em consideração aspectos como multietnicidade, pluralidade e diversidade, educação e conhecimentos \textbf{indígenas}, autodeterminação, comunidade educativa \textbf{indígena}, bem como os princípios dessa educação.
Atualmente o \textbf{Estado} de Rondônia atende 109 escolas \textbf{indígenas}, com um total 3.468 estudantes, acompanhados por 14 Coordenações de Educação Escolares \textbf{Indígenas} Regionais, localizadas nas Coordenadorias Regionais de Educação.
Ao todo são atendidas 54 etnias falantes de 29 ( vinte e nove ) línguas \textbf{indígenas}, e 3 ( três ) dialetos onde aparecem comunidades que vão desde agrupamentos humanos fragmentados de menos de uma dezena de indivíduos, como os Karipuna, até comunidades de mais de mil indivíduos como os povos Oro Wari.
As escolas estão localizadas em áreas \textbf{indígenas}, em sua maioria de difícil acesso, sendo terrestres ou fluviais.
Das 109 escolas \textbf{indígenas} que constam no Quadro abaixo, 3 ( três ) escolas estão administrativamente sob a responsabilidade do Município de Chupinguaia, sendo que o \textbf{Estado} fornece professores e coordenação escolar, todas as demais são atendidas integralmente pelo \textbf{Estado}.
A Educação Escolar \textbf{Indígena} é um direito assegurado aos povos \textbf{indígenas} e seu ensino é pautado numa educação diferenciada.
Para que isso aconteça, a Legislação Escolar \textbf{Indígena} ( 1998 ) se pauta em duas vertentes que devem ser sempre convergentes entre si : a de propiciar acesso ao conhecimento dito universal ; e a de ensejar práticas escolares que permitam o respeito e a sistematização de saberes e conhecimentos tradicionais.
Ambas as vertentes hoje estão em construção ; no entanto, podemos observar que há um aparato legal que já ampara essa educação e que respalda um ensino específico, multicultural e próprio aos povos \textbf{indígenas}.
Conforme pode ser observado no quadro abaixo : EDUCAÇÃO ESCOLAR \textbf{INDÍGENA} | TEMÁTICA | CONTEÚDO | |---|---| | Multietnicidade, Pluralidade e Diversidade | Reconhecer e manter situação de comunicação entre as experiências socioculturais, linguísticas e históricas. | | Bilinguismo, Multilinguismo e Plurilinguismo | As \textbf{populações} \textbf{Indígenas} e a convivência com várias situações linguísticas : monolíngues tem apenas uma Língua em seu repertório, Os \textbf{Bilíngues} têm duas, os Trilíngues têm três ; O ensino \textbf{Bilíngue} em escolas \textbf{indígenas} ; O Bilinguismo Competente fala e escreve ; O Bilinguismo Mínimo – falante sem escrita ; E os povos Multilíngues são aqueles que convivem, com uma diversidade linguística, ocasionando um vasto repertório verbal, realidade ocorre especialmente em regiões fronteiriças. | | Pluralidade Cultural | Fortalecer sua \textbf{Identidade} Étnica, suas Tradições e práticas Culturais ; Tornar o espaço escolar possível de interculturalidade ; Contribuir para o projeto de autonomia dos Povos \textbf{Indígenas}, a partir de seus Projetos Históricos, desenvolvendo novas estratégias de sobrevivência, linguística e Cultural ; A diversidade das Sociedades \textbf{indígenas} ; Terra, com Preservação da Biodiversidade, Auto sustentação.
Conservação da \textbf{Fauna} ; Evitar perda da biodiversidade da \textbf{Fauna} é atribuído à devastação dos habitats naturais e agravada por práticas como a caça e o tráfico de animais ; Conservação da Flora como Equilíbrio ; Como os povos \textbf{indígenas} e quilombolas ajudam a ampliar a diversidade da \textbf{fauna} e da \textbf{flora} local porque têm formas únicas de viver e ocupar um lugar ; Preservar Sementes Tradicionais ; Incentivar a auto – sustentação junto às Comunidades \textbf{Indígenas}, visando o enfrentamento dos problemas Ambientais oriundos da convivência com a sociedade envolvente ; Uma relação de respeito com a água, que é passada de geração para geração desde os ancestrais ; A água é um ser vivo, divino de uso e propriedade comunitária que deve ser compartilhada ; Mitos e explicações culturais de cada povo Mitos \textbf{indígenas} são narrativos de tradição oral, passadas de geração em geração, que funcionam como um arquivo de conhecimento ; os mitos \textbf{indígenas}, contam o jeito certo de viver em sociedade, de fazer as festas e os rituais, de fazer roça, de caçar, de pescar, e fazer rede, cestos, dentre tantas outras coisas importantes para a vida ; O Mito é um método pedagógico de narrar os acontecimentos dos ancestrais Direitos \textbf{Indígenas} na Sociedade nacional ; Respeito, Preservação a integridade Física e Moral dos povos \textbf{Indígenas} ; Exercício dos direitos nas próprias Comunidades ; Conhecer a história dos grandes movimentos \textbf{indígenas} no Brasil ; Trazer para a escola a discussão da Ética, recolocar valores Particulares às sociedades \textbf{Indígenas} ; Ser um instrumento de interlocução entre os Saberes da Sociedade \textbf{Indígena} e a aquisição de outros conhecimentos não \textbf{indígenas} ; Direito a Construção de uma Proposta Curricular contextualizada Quilombola ; Direito a lutas por demarcação territorial de sua \textbf{ancestralidade} Quilombola ; Direito a oferta do ensino Étnico- \textbf{Racial} quilombola ; Direito a manutenção de seus usos, costumes e tradições relacionada a resistência à opressão histórica ; |

% matched lemmas: ancestralidade, bilíngue, estado, fauna, flora, identidade, indígena, população, racial
\end{textsample}
